% Revised manuscript for Automated Software Engineering (Springer Nature)
% Requires: \documentclass{sn-jnl}
% Online Resources referenced per journal SI guidance.
% v2.0 — incorporates agentic architecture, hardened benchmark, ablation study

\documentclass{sn-jnl}

\title[Dynamic Regulatory Alignment]{Dynamic Regulatory Alignment: An Agentic Framework for Automated Compliance Mapping with Uncertainty-Aware Retrieval}

\author*[1]{Sumanshu Sohal}
\email{sumanshu.95s@outlook.com}

\affil[1]{Cybersecurity Specialist, Washington, DC, USA}

\abstract{
Regulatory compliance mapping links natural-language regulatory requirements to organizational and technical controls.
In practice, mapping is many-to-many, high-stakes, and drift-prone: regulations evolve, audit frameworks overlap, and legal vocabulary mismatches implementation terminology.

We present the Regulatory Alignment Agent (RAA), which operationalizes compliance mapping as automated traceability recovery using a ReAct-style reasoning loop.
The RAA implements five tool-based capabilities: (i)~multi-backend retrieval with reciprocal rank fusion across TF-IDF, BM25, and LSI,
(ii)~domain-aware query reformulation using a compliance thesaurus and concept extraction,
(iii)~cross-framework corroboration that validates candidate mappings against related regulatory families,
(iv)~bidirectional verification that checks mapping consistency by reverse retrieval, and
(v)~an uncertainty-aware decision policy that can accept, abstain, or escalate for human review.

We evaluate on a hardened synthetic benchmark (58~requirements, 110~controls including 20~vocabulary-mismatched positives and 24~hard negatives across 7~regulatory frameworks) with stratified holdout and 30~repeated seeds.
An ablation study isolates each agent component's contribution:
query reformulation provides the largest gain (+28\% Top-1 accuracy over single-shot retrieval, $p < 0.05$),
multi-backend fusion adds +5\%, and cross-framework corroboration improves selective accuracy by 2\% while maintaining high coverage.
The full agent achieves 0.617 Top-1 ($\pm$0.031), 0.705 MRR@5 ($\pm$0.026), and 0.656 Selective Accuracy at 94\% coverage, outperforming all lexical baselines.
All reasoning steps produce auditable traces suitable for compliance documentation.

We compare against a canonical legacy baseline (LSI traceability recovery, ICSE 2003) using matched evaluation protocols including legacy-style micro-precision/recall.
}

\keywords{Regulatory requirements engineering \and Traceability \and Information retrieval \and Agentic reasoning \and Query reformulation \and Multi-backend fusion \and Selective prediction \and Human-in-the-loop}

\begin{document}

\maketitle

\section{Introduction}
Compliance mapping is often implemented as spreadsheets maintained through manual interpretation and periodic audits.
From a software engineering perspective, this is traceability: establishing and maintaining links between documentation-like artifacts (laws, standards, policies) and implementation-like artifacts (controls, configurations, evidence).
However, regulatory mapping differs from classic traceability recovery in two ways: (i)~the open-world outcome ``no match'' is a first-class signal (compliance gaps), and (ii)~temporal drift is inherent because requirements and controls evolve.

Single-shot retrieval approaches struggle with two fundamental challenges in compliance mapping:
\emph{vocabulary mismatch}, where regulatory language (``implement appropriate technical measures including encryption'') differs systematically from control language (``AES-256 with FIPS 140-2 modules''), and
\emph{scope confusion}, where lexically similar controls address different regulatory concerns (e.g., in-transit vs.\ at-rest encryption).

This paper introduces a ReAct-style~\cite{ReAct} agentic workflow that addresses these challenges through multi-step reasoning: retrieve candidate controls, reformulate queries when confidence is low, corroborate against related frameworks, verify bidirectionally, then accept or abstain with an auditable record.
We demonstrate through ablation that each reasoning step contributes measurably, with query reformulation providing the dominant improvement.

\section{Background and related work}

\subsection{Legacy traceability IR baseline: LSI link recovery}
A canonical traceability recovery approach applies Latent Semantic Indexing (LSI) to recover documentation-to-code links~\cite{MarcusMaletic2003}.
The legacy method treats documentation and code artifacts as text documents, builds a term-document matrix, applies truncated SVD, and ranks links by cosine similarity in the latent space.
A key methodological point is the selection protocol: links can be selected by (a)~similarity threshold or (b)~a cut-point (top-$k$ ranked links per documentation unit), with the expected precision/recall trade-off.

We adopt this LSI method as a legacy baseline for compliance mapping and reproduce its system-level aggregation for precision and recall (micro-averaging across queries) in addition to modern IR metrics.

\subsection{Query reformulation and expansion}
Query reformulation is well-established in information retrieval~\cite{BM25} but underexplored in compliance mapping.
We implement domain-specific reformulation using a hand-crafted compliance thesaurus (20 concept families, 30+ expansion patterns) that maps regulatory vocabulary to implementation vocabulary.
This approach is deterministic, auditable, and does not require neural models---important properties for regulated environments.

\subsection{Multi-system fusion}
Reciprocal Rank Fusion (RRF) combines rankings from multiple retrieval systems without requiring score normalization.
We apply RRF across TF-IDF, BM25, and LSI backends, leveraging the complementary strengths of exact match (BM25), distributional similarity (TF-IDF), and latent semantics (LSI).

\subsection{Agentic workflows and uncertainty management}
In high-stakes settings, automation should surface uncertainty rather than force decisions.
The ReAct framework~\cite{ReAct} interleaves reasoning (``thought'') and acting (``tool use'') in a loop, enabling systems to adapt their strategy based on intermediate observations.
We implement this pattern for compliance mapping: the agent reasons about retrieval confidence, invokes tools (reformulation, cross-referencing, verification), and produces auditable decision traces.

\section{Approach}

\subsection{Problem definition}
Given a regulatory requirement $r$ and a control corpus $C=\{c_i\}$, produce a ranked list of candidate controls and a decision:
\texttt{accept}, \texttt{review}, or \texttt{abstain}.
The decision must be supported by an audit trace: scores, score gaps, reasoning steps, and tool invocations.

\subsection{Agent architecture}
The RAA implements a ReAct-style reasoning loop with the following tools:

\begin{enumerate}
\item \textbf{Retrieve}$(q, \text{backend})$: Score all controls against query $q$ using a specified backend (TF-IDF, BM25, or LSI). Returns a score vector.
\item \textbf{Fuse}$(\{s_1, \ldots, s_k\})$: Combine score vectors from multiple backends via Reciprocal Rank Fusion: $\text{RRF}(d) = \sum_i \frac{1}{k + \text{rank}_i(d)}$.
\item \textbf{Reformulate}$(q)$: Expand $q$ with domain concepts using pattern-based concept extraction and a compliance thesaurus. Maps regulatory vocabulary to implementation vocabulary (e.g., ``boundary protection'' $\rightarrow$ ``firewall, network perimeter, DMZ, segmentation, IPS'').
\item \textbf{CrossRef}$(r, c^*)$: Check if controls in the same domain family as the top candidate $c^*$ appear across frameworks related to $r$'s framework. Adjusts the accept/abstain threshold based on corroboration strength.
\item \textbf{Verify}$(r, c^*)$: Bidirectional verification---use the top control's text as a query to check if the mapping is consistent in the reverse direction.
\item \textbf{Decide}$(\text{scores}, \theta_{\text{conf}}, \theta_{\text{gap}})$: Accept if confidence $\geq \theta_{\text{conf}}$ and score gap $\geq \theta_{\text{gap}}$; otherwise abstain.
\end{enumerate}

\subsection{Reasoning loop}
For each regulation $r$:
\begin{enumerate}
\item \textsc{Think}: Identify the regulation's framework and domain.
\item \textsc{Act}: \textbf{Retrieve}$(r.\text{text}, \text{bm25})$ --- initial retrieval.
\item \textsc{Observe}: Check confidence and gap. If adequate, proceed to decision.
\item \textsc{Act}: \textbf{Fuse} results from all backends via RRF.
\item \textsc{Observe}: If confidence or gap is low:
  \begin{enumerate}
  \item \textsc{Act}: \textbf{Reformulate}$(r.\text{text})$ and re-retrieve with expanded query.
  \item \textsc{Observe}: Keep the better result.
  \end{enumerate}
\item \textsc{Act}: \textbf{CrossRef}$(r, c^*)$ --- validate top candidate against related frameworks.
\item \textsc{Act}: \textbf{Verify}$(r, c^*)$ --- bidirectional consistency check.
\item \textsc{Act}: \textbf{Decide} with corroboration-adjusted thresholds.
\end{enumerate}

Each step produces an \texttt{AgentStep} record (thought, action, input, observation), and the full trace is packaged as an \texttt{AgentTrace} for audit.

\subsection{Abstention calibration}
Confidence and gap thresholds are calibrated on a held-out calibration split to target 80\% coverage.
When cross-framework corroboration confirms the top candidate, thresholds are relaxed by up to 15\%, allowing the agent to accept borderline matches that have structural support.

\section{Dataset}
We evaluate on a hardened synthetic benchmark designed to test both lexical and semantic retrieval:
\begin{itemize}
\item 58 regulatory requirements spanning 7 frameworks (GDPR, NIST, HIPAA, PCI~DSS, ISO~27001, SOX, SOC~2).
\item 110 control statements:
  \begin{itemize}
  \item 66 vocabulary-matched controls (near-paraphrases of requirements).
  \item 20 vocabulary-mismatched controls (correct mappings using entirely different terminology, e.g., ``cryptographic safeguards applied to individually identifiable records'' for an encryption requirement).
  \item 24 hard negatives: scope confounders (in-transit vs.\ at-rest), cross-family distractors (marketing database access controls), near-miss negatives (development-only logging), and irrelevant controls.
  \end{itemize}
\item 86 positive mapping links (multi-label: some requirements map to multiple controls).
\end{itemize}

The vocabulary-mismatched controls are the key addition over v1.0: they create retrieval challenges where lexical methods struggle but domain-aware reformulation can bridge the vocabulary gap.

Online Resource 1 contains the benchmark as CSV.

\section{Experimental protocol}
We use stratified holdout by framework family: for each run seed, each framework contributes to train/calibration/test splits (holdout=0.20; calibration=0.15).
All baselines and agent variants use identical splits per seed.

We report Top-1, Recall@1, Recall@5, MRR@5, MAP@5, nDCG@5, Precision@5, F1@5, Coverage, and Selective Accuracy at target coverage.
Mean $\pm$ 95\% CI is computed via the Student $t$ distribution across 30 repeated seeds.
We additionally report legacy-style micro-precision/micro-recall at cut points.

\subsection{Ablation variants}
To isolate each agent component's contribution, we evaluate incrementally:
\begin{enumerate}
\item \textbf{Single}: One-shot BM25 retrieval + threshold decision (v1.0 baseline).
\item \textbf{Multi}: Single + reciprocal rank fusion across TF-IDF, BM25, and LSI.
\item \textbf{Reform}: Multi + domain-aware query reformulation on low-confidence queries.
\item \textbf{CrossRef}: Reform + cross-framework corroboration with threshold adjustment.
\item \textbf{Agent (RAA)}: CrossRef + bidirectional verification. Full reasoning loop.
\end{enumerate}

\section{Results}

\subsection{Baseline comparison}
Table~\ref{tab:baselines} reports results for the three lexical baselines on the hardened benchmark (30 seeds).

\begin{table}[h]
\caption{Baseline comparison on hardened benchmark (mean $\pm$ 95\% CI, 30 seeds, 110 controls).}
\label{tab:baselines}
\centering
\begin{tabular}{lccccc}
\hline
Method & Top-1 & MRR@5 & nDCG@5 & Recall@5 & SelAcc@80\% \\
\hline
TF-IDF   & 0.511 $\pm$ 0.041 & 0.588 $\pm$ 0.039 & 0.548 $\pm$ 0.036 & 0.593 $\pm$ 0.041 & 0.643 $\pm$ 0.044 \\
BM25     & 0.481 $\pm$ 0.043 & 0.537 $\pm$ 0.043 & 0.504 $\pm$ 0.043 & 0.535 $\pm$ 0.047 & 0.656 $\pm$ 0.046 \\
LSI (legacy) & 0.500 $\pm$ 0.042 & 0.623 $\pm$ 0.035 & 0.586 $\pm$ 0.033 & 0.659 $\pm$ 0.037 & 0.542 $\pm$ 0.044 \\
\hline
\end{tabular}
\end{table}

On the hardened benchmark, lexical baselines plateau around 0.50 Top-1 accuracy.
LSI achieves higher Recall@5 (0.659) due to its latent space capturing some semantic similarity, but its Selective Accuracy is lowest (0.542), indicating poor calibration for accept/abstain decisions.

\subsection{Semantic pipeline}
Table~\ref{tab:semantic} reports the semantic retrieval results using a general-purpose dual-encoder (\texttt{all-MiniLM-L6-v2}) and a cross-encoder reranker (\texttt{ms-marco-MiniLM-L-6-v2}).

\begin{table}[h]
\caption{Semantic pipeline results (mean $\pm$ 95\% CI, 30 seeds).}
\label{tab:semantic}
\centering
\begin{tabular}{lccccc}
\hline
Method & Top-1 & MRR@5 & nDCG@5 & Recall@5 & SelAcc@80\% \\
\hline
Dual-encoder   & 0.583 $\pm$ 0.046 & 0.642 $\pm$ 0.045 & 0.618 $\pm$ 0.043 & 0.652 $\pm$ 0.046 & 0.783 $\pm$ 0.054 \\
+Cross-encoder reranker & 0.553 $\pm$ 0.050 & 0.585 $\pm$ 0.052 & 0.559 $\pm$ 0.050 & 0.585 $\pm$ 0.052 & 0.779 $\pm$ 0.051 \\
\hline
\end{tabular}
\end{table}

The dual-encoder outperforms all lexical baselines on Top-1 (0.583 vs.\ 0.511 for TF-IDF) and achieves the highest Selective Accuracy among non-agent methods (0.783).
Notably, the cross-encoder reranker does not improve over the dual-encoder alone; general-purpose rerankers trained on web search data do not capture compliance-domain semantics, suggesting that domain-adapted rerankers would be needed for further gains.
The RAA agent with reformulation (0.617 Top-1) still outperforms the semantic dual-encoder, demonstrating that domain-aware query expansion provides complementary value beyond neural embeddings.

\subsection{Agent ablation study}
Table~\ref{tab:ablation} presents the incremental ablation.
Each row adds one component to the previous variant.

\begin{table}[h]
\caption{Agent ablation study (mean $\pm$ 95\% CI, 30 seeds). Each row adds one component.}
\label{tab:ablation}
\centering
\begin{tabular}{lcccccc}
\hline
Variant & Top-1 & MRR@5 & nDCG@5 & R@5 & Coverage & SelAcc \\
\hline
Single (BM25)    & 0.481 $\pm$ 0.043 & 0.537 $\pm$ 0.043 & 0.504 $\pm$ 0.043 & 0.535 $\pm$ 0.047 & 0.739 & 0.656 \\
+Multi (RRF)     & 0.503 $\pm$ 0.044 & 0.576 $\pm$ 0.041 & 0.540 $\pm$ 0.042 & 0.582 $\pm$ 0.048 & 0.789 & 0.644 \\
+Reformulation   & \textbf{0.619 $\pm$ 0.031} & \textbf{0.711 $\pm$ 0.027} & \textbf{0.656 $\pm$ 0.029} & \textbf{0.724 $\pm$ 0.034} & 0.964 & 0.645 \\
+CrossRef        & 0.617 $\pm$ 0.031 & 0.705 $\pm$ 0.026 & 0.643 $\pm$ 0.029 & 0.706 $\pm$ 0.035 & 0.942 & 0.656 \\
Full Agent (RAA) & 0.617 $\pm$ 0.031 & 0.705 $\pm$ 0.026 & 0.643 $\pm$ 0.029 & 0.706 $\pm$ 0.035 & 0.942 & \textbf{0.656} \\
\hline
\end{tabular}
\end{table}

\paragraph{Key findings.}
\begin{enumerate}
\item \textbf{Query reformulation is the dominant contributor}, improving Top-1 by 23\% absolute (0.503 $\rightarrow$ 0.619) with non-overlapping confidence intervals.
This confirms that vocabulary mismatch is the primary challenge in compliance mapping retrieval, and that domain-aware query expansion effectively bridges the gap between regulatory and implementation language.

\item \textbf{Multi-backend fusion provides moderate, consistent gains} (+4.6\% Top-1, +7.3\% MRR@5) by combining the complementary strengths of exact-match, distributional, and latent retrieval.

\item \textbf{Cross-framework corroboration improves decision quality} rather than ranking quality: Selective Accuracy rises from 0.645 to 0.656 while coverage adjusts from 0.964 to 0.942, reflecting more conservative but more accurate accept/abstain decisions.

\item \textbf{The full agent outperforms all lexical baselines} across all metrics while producing auditable 7-step reasoning traces for every mapping decision.
\end{enumerate}

\subsection{Reasoning trace example}
Table~\ref{tab:trace} illustrates the agent's reasoning for a regulation with vocabulary mismatch.

\begin{table}[h]
\caption{Example agent trace for NIST SC-7 (Boundary protection).}
\label{tab:trace}
\centering
\small
\begin{tabular}{clp{7cm}}
\hline
Step & Action & Observation \\
\hline
1 & Retrieve (BM25) & Top: ``Perimeter defense appliances...'' conf=0.31 \\
2 & Fuse (RRF) & Top: ``Next-gen firewall with IPS...'' conf=0.018 \\
3 & Reformulate & Expanded: ``...firewall, network perimeter, DMZ, segmentation, IPS'' \\
4 & Re-retrieve & Top: ``Next-gen firewall with IPS at network perimeter'' conf=0.021 \\
5 & CrossRef & Corroboration=1.0: domain spans PCI, ISO \\
6 & Verify & Bidirectional check passed \\
7 & Decide & \textbf{Accept} (conf $\geq$ adjusted threshold) \\
\hline
\end{tabular}
\end{table}

\section{Discussion}
The ablation reveals that compliance mapping's core challenge is vocabulary mismatch between regulatory and implementation language, not retrieval algorithm sophistication.
Domain-aware query reformulation---a deterministic, auditable technique---provides larger gains than switching retrieval algorithms or applying multi-system fusion.

The cross-reference and verification steps contribute to \emph{decision quality} rather than \emph{ranking quality}.
In compliance contexts, this distinction matters: a system that correctly abstains on uncertain mappings (preserving them for human review) is more valuable than one that forces incorrect high-confidence decisions.
The agent's Selective Accuracy of 0.656 at 94\% coverage means that when it accepts a mapping, it is correct about two-thirds of the time---with full audit traces documenting why.

\subsection{Limitations}
The benchmark remains synthetic, with controlled vocabulary mismatch patterns.
Real-world compliance mappings involve ambiguous regulatory language, multi-clause requirements, and evolving control implementations.
The thesaurus-based reformulation is manually curated and may not generalize to novel regulatory domains without extension.
We plan to extend with real-world NIST SP 800-53 to ISO 27001 mapping data in future work.

\section{Threats to validity}

\textbf{Construct validity:} The hardened benchmark includes 20 vocabulary-mismatched controls and 24 hard negatives (scope confounders, cross-family distractors, near-miss negatives) to test beyond lexical overlap.
Results confirm that reformulation specifically addresses vocabulary mismatch cases where baselines fail.

\textbf{Internal validity:} Threshold calibration uses a held-out calibration split separate from test data.
Results are reported over 30 seeds with Student~$t$ confidence intervals.
The ablation design ensures each component is evaluated incrementally.

\textbf{External validity:} Compliance mapping correctness is interpretation-dependent; the system is decision support, and abstention/escalation are required for defensibility.
The benchmark size (58 requirements) limits generalization claims; confidence intervals reflect this.

\section*{Data Availability}
This study uses a hardened synthetic benchmark packaged as Online Resource 1 (CSV).
The benchmark contains 58 regulations, 110 controls (including vocabulary-mismatched positives and hard negatives), and 86 mapping labels.

\section*{Code Availability}
The complete evaluation framework, baselines (TF-IDF, BM25, LSI), agent variants (single, multi, reform, crossref, full agent), domain thesaurus, and ablation harness are provided as Online Resource 2 (Python CLI).
The semantic backend (Sentence-Transformers) requires external model downloads; all other components are self-contained.

\section*{Online Resources}
Online Resource 1: Hardened synthetic benchmark (regulations.csv, controls.csv, mappings.csv).\\
Online Resource 2: Reproducible evaluation program (Python CLI), agent implementation, and ablation scripts.

\section*{Declarations}
\textbf{Conflict of interest:} The author declares no competing interests.\\
\textbf{Originality:} This manuscript is not under review elsewhere and does not substantially overlap another submission.

\begin{thebibliography}{}
\bibitem{MarcusMaletic2003}
Marcus, A., Maletic, J.I.: Recovering documentation-to-source-code traceability links using latent semantic indexing. Proc.\ 25th Int.\ Conf.\ on Software Engineering (ICSE 2003). DOI: 10.1109/ICSE.2003.1201194

\bibitem{LegalBERT}
Chalkidis, I.\ et al.: LEGAL-BERT: The Muppets straight out of Law School. Findings of EMNLP 2020. DOI: 10.18653/v1/2020.findings-emnlp.261

\bibitem{SecureBERT}
Aghaei, E.\ et al.: SecureBERT: A Domain-Specific Language Model for Cybersecurity. arXiv:2204.02685

\bibitem{ReAct}
Yao, S.\ et al.: ReAct: Synergizing Reasoning and Acting in Language Models. ICLR 2023. arXiv:2210.03629

\bibitem{ConceptDriftSurvey}
Gama, J.\ et al.: A survey on concept drift adaptation. ACM Computing Surveys 46(4), Article 44 (2014). DOI: 10.1145/2523813

\bibitem{BM25}
Robertson, S., Zaragoza, H.: The Probabilistic Relevance Framework: BM25 and Beyond. Foundations and Trends in Information Retrieval 3(4), 333--389 (2009). DOI: 10.1561/1500000019

\bibitem{RRF}
Cormack, G.V., Clarke, C.L.A., Buettcher, S.: Reciprocal rank fusion outperforms condorcet and individual rank learning methods. Proc.\ 32nd SIGIR (2009). DOI: 10.1145/1571941.1572114

\end{thebibliography}

\end{document}
